% Algebra 2 Lesson 2-4 fast-track packet: Complex Numbers and Operations (enVision A2, Example 3 only; brief lesson)
% Shared preamble for the "fast" (hacked-down) Algebra 2 lesson packets, Oct 2026.
% Each packet: % Shared preamble for the "fast" (hacked-down) Algebra 2 lesson packets, Oct 2026.
% Each packet: % Shared preamble for the "fast" (hacked-down) Algebra 2 lesson packets, Oct 2026.
% Each packet: \input{fast-packet-template} then \begin{document} ... \end{document}.
\documentclass[12pt]{article}
\usepackage[margin=0.7in]{geometry}
\usepackage{amsmath,amssymb}
\usepackage{booktabs,tabularx,array}
\usepackage{enumitem}
\usepackage[T1]{fontenc}
\usepackage[utf8]{inputenc}
\usepackage{lmodern}
\usepackage{tikz}
\usepackage{pgfplots}
\pgfplotsset{compat=1.18}
\usepackage{xcolor}
\setlist{nosep,leftmargin=1.4em}
\setlength{\parindent}{0pt}
\setlength{\parskip}{6pt}
\renewcommand{\arraystretch}{1.35}
\newcolumntype{Y}{>{\raggedright\arraybackslash}X}
\pagestyle{empty}

% blanks and workspace
\newcommand{\blank}[1][2.6cm]{\rule[-2pt]{#1}{0.5pt}}
\newcommand{\lines}[1][1]{\par\vspace{2pt}\foreach \i in {1,...,#1}{\rule{0pt}{22pt}\rule{\linewidth}{0.4pt}\par}}
\newcommand{\wspart}[2]{\vspace{4pt}{\large\bfseries #1: #2}\par\vspace{-6pt}\rule{\textwidth}{0.4pt}\par}

% header: \packethead{Lesson 1-3}{Piecewise-Defined Functions}
\newcommand{\packethead}[2]{%
{\LARGE\bfseries #2}\hfill Name \blank[5cm]\\[2pt]
Algebra 2 \quad #1 (fast track) \hfill Date \blank[2.5cm] \quad Period \blank[1.2cm]\par\vspace{2pt}}

% boxed goals: \goals{learning goal}{language goal}
\newcommand{\goals}[2]{%
\fbox{\parbox{\dimexpr\textwidth-2\fboxsep-2\fboxrule}{%
\textbf{Learning goal:} #1\\[2pt]
\textbf{Language goal:} #2}}\par}

% quiz target box: \quiztarget{text}
\newcommand{\quiztarget}[1]{%
\vspace{4pt}\colorbox{gray!12}{\parbox{\dimexpr\textwidth-2\fboxsep}{\textbf{On the quiz:} #1}}\par}

% IXL table row helper: code & skill & which quiz question it feeds
\newcommand{\ixlrow}[3]{\texttt{#1} & #2 & #3 \\}

% one grid style for every graph; override xmin/xmax/... per plot
\pgfplotsset{sheetgrid/.style={
  axis lines=middle, grid=both, grid style={gray!45},
  xmin=-10.5, xmax=10.5, ymin=-10.5, ymax=10.5,
  xtick={-10,-8,...,10}, ytick={-10,-8,...,10},
  x=0.28cm, y=0.28cm,
  ticklabel style={font=\scriptsize, fill=white, inner sep=1pt},
  xlabel={$x$}, ylabel={$y$},
  axis line style={thick}, clip=true}}
% unit grid for small windows
\pgfplotsset{unitgrid/.style={sheetgrid,
  xmin=-5.5, xmax=5.5, ymin=-5.5, ymax=5.5,
  xtick={-5,...,5}, ytick={-5,...,5}, x=0.5cm, y=0.5cm}}
% open and closed endpoints
\tikzset{closed/.style={circle, fill, inner sep=1.6pt},
         open/.style={circle, draw, fill=white, inner sep=1.4pt}}
 then \begin{document} ... \end{document}.
\documentclass[12pt]{article}
\usepackage[margin=0.7in]{geometry}
\usepackage{amsmath,amssymb}
\usepackage{booktabs,tabularx,array}
\usepackage{enumitem}
\usepackage[T1]{fontenc}
\usepackage[utf8]{inputenc}
\usepackage{lmodern}
\usepackage{tikz}
\usepackage{pgfplots}
\pgfplotsset{compat=1.18}
\usepackage{xcolor}
\setlist{nosep,leftmargin=1.4em}
\setlength{\parindent}{0pt}
\setlength{\parskip}{6pt}
\renewcommand{\arraystretch}{1.35}
\newcolumntype{Y}{>{\raggedright\arraybackslash}X}
\pagestyle{empty}

% blanks and workspace
\newcommand{\blank}[1][2.6cm]{\rule[-2pt]{#1}{0.5pt}}
\newcommand{\lines}[1][1]{\par\vspace{2pt}\foreach \i in {1,...,#1}{\rule{0pt}{22pt}\rule{\linewidth}{0.4pt}\par}}
\newcommand{\wspart}[2]{\vspace{4pt}{\large\bfseries #1: #2}\par\vspace{-6pt}\rule{\textwidth}{0.4pt}\par}

% header: \packethead{Lesson 1-3}{Piecewise-Defined Functions}
\newcommand{\packethead}[2]{%
{\LARGE\bfseries #2}\hfill Name \blank[5cm]\\[2pt]
Algebra 2 \quad #1 (fast track) \hfill Date \blank[2.5cm] \quad Period \blank[1.2cm]\par\vspace{2pt}}

% boxed goals: \goals{learning goal}{language goal}
\newcommand{\goals}[2]{%
\fbox{\parbox{\dimexpr\textwidth-2\fboxsep-2\fboxrule}{%
\textbf{Learning goal:} #1\\[2pt]
\textbf{Language goal:} #2}}\par}

% quiz target box: \quiztarget{text}
\newcommand{\quiztarget}[1]{%
\vspace{4pt}\colorbox{gray!12}{\parbox{\dimexpr\textwidth-2\fboxsep}{\textbf{On the quiz:} #1}}\par}

% IXL table row helper: code & skill & which quiz question it feeds
\newcommand{\ixlrow}[3]{\texttt{#1} & #2 & #3 \\}

% one grid style for every graph; override xmin/xmax/... per plot
\pgfplotsset{sheetgrid/.style={
  axis lines=middle, grid=both, grid style={gray!45},
  xmin=-10.5, xmax=10.5, ymin=-10.5, ymax=10.5,
  xtick={-10,-8,...,10}, ytick={-10,-8,...,10},
  x=0.28cm, y=0.28cm,
  ticklabel style={font=\scriptsize, fill=white, inner sep=1pt},
  xlabel={$x$}, ylabel={$y$},
  axis line style={thick}, clip=true}}
% unit grid for small windows
\pgfplotsset{unitgrid/.style={sheetgrid,
  xmin=-5.5, xmax=5.5, ymin=-5.5, ymax=5.5,
  xtick={-5,...,5}, ytick={-5,...,5}, x=0.5cm, y=0.5cm}}
% open and closed endpoints
\tikzset{closed/.style={circle, fill, inner sep=1.6pt},
         open/.style={circle, draw, fill=white, inner sep=1.4pt}}
 then \begin{document} ... \end{document}.
\documentclass[12pt]{article}
\usepackage[margin=0.7in]{geometry}
\usepackage{amsmath,amssymb}
\usepackage{booktabs,tabularx,array}
\usepackage{enumitem}
\usepackage[T1]{fontenc}
\usepackage[utf8]{inputenc}
\usepackage{lmodern}
\usepackage{tikz}
\usepackage{pgfplots}
\pgfplotsset{compat=1.18}
\usepackage{xcolor}
\setlist{nosep,leftmargin=1.4em}
\setlength{\parindent}{0pt}
\setlength{\parskip}{6pt}
\renewcommand{\arraystretch}{1.35}
\newcolumntype{Y}{>{\raggedright\arraybackslash}X}
\pagestyle{empty}

% blanks and workspace
\newcommand{\blank}[1][2.6cm]{\rule[-2pt]{#1}{0.5pt}}
\newcommand{\lines}[1][1]{\par\vspace{2pt}\foreach \i in {1,...,#1}{\rule{0pt}{22pt}\rule{\linewidth}{0.4pt}\par}}
\newcommand{\wspart}[2]{\vspace{4pt}{\large\bfseries #1: #2}\par\vspace{-6pt}\rule{\textwidth}{0.4pt}\par}

% header: \packethead{Lesson 1-3}{Piecewise-Defined Functions}
\newcommand{\packethead}[2]{%
{\LARGE\bfseries #2}\hfill Name \blank[5cm]\\[2pt]
Algebra 2 \quad #1 (fast track) \hfill Date \blank[2.5cm] \quad Period \blank[1.2cm]\par\vspace{2pt}}

% boxed goals: \goals{learning goal}{language goal}
\newcommand{\goals}[2]{%
\fbox{\parbox{\dimexpr\textwidth-2\fboxsep-2\fboxrule}{%
\textbf{Learning goal:} #1\\[2pt]
\textbf{Language goal:} #2}}\par}

% quiz target box: \quiztarget{text}
\newcommand{\quiztarget}[1]{%
\vspace{4pt}\colorbox{gray!12}{\parbox{\dimexpr\textwidth-2\fboxsep}{\textbf{On the quiz:} #1}}\par}

% IXL table row helper: code & skill & which quiz question it feeds
\newcommand{\ixlrow}[3]{\texttt{#1} & #2 & #3 \\}

% one grid style for every graph; override xmin/xmax/... per plot
\pgfplotsset{sheetgrid/.style={
  axis lines=middle, grid=both, grid style={gray!45},
  xmin=-10.5, xmax=10.5, ymin=-10.5, ymax=10.5,
  xtick={-10,-8,...,10}, ytick={-10,-8,...,10},
  x=0.28cm, y=0.28cm,
  ticklabel style={font=\scriptsize, fill=white, inner sep=1pt},
  xlabel={$x$}, ylabel={$y$},
  axis line style={thick}, clip=true}}
% unit grid for small windows
\pgfplotsset{unitgrid/.style={sheetgrid,
  xmin=-5.5, xmax=5.5, ymin=-5.5, ymax=5.5,
  xtick={-5,...,5}, ytick={-5,...,5}, x=0.5cm, y=0.5cm}}
% open and closed endpoints
\tikzset{closed/.style={circle, fill, inner sep=1.6pt},
         open/.style={circle, draw, fill=white, inner sep=1.4pt}}

\begin{document}

\packethead{Lesson 2-4}{Complex Numbers and Operations}
\goals{I can multiply complex numbers and write each product in the form $a+bi$.}{I can explain each step of a product using the words \emph{imaginary unit}, \emph{real part}, \emph{imaginary part}, and \emph{complex conjugates}.}

\quiztarget{(Savvas numbers) (1) Solve $x^2=-64$ over the complex numbers (full Savvas quiz only). (2) Bonus: subtract two complex numbers; combine the real parts, then the imaginary parts (Practice b). (3) Multiply a complex number by its \emph{complex conjugate}: the $i$ terms cancel and the product is a real number (Example 3B, Try It 3b). (4) Multiply three factors (an imaginary number times two binomials) and write the answer as $a+bi$ (Example 3A and 3B, Practice a). (5) Match products with their complex factors (full Savvas quiz only).}

\fbox{\parbox{\dimexpr\textwidth-2\fboxsep-2\fboxrule}{\small
\textbf{Complex numbers.} The \textbf{imaginary unit} $i$ is the principal square root of $-1$. Then $i^2=-1$.
A \textbf{complex number} can be written in the form $a+bi$, where $a$ and $b$ are real numbers: $a$ is the \textbf{real part}, $b$ is the \textbf{imaginary part}, and $bi$ is the imaginary term. For example: $-6+4i$ ($a=-6$, $b=4$), \ $7-i\sqrt2$ ($a=7$, $b=-\sqrt2$), \ $0.5i$ ($a=0$, $b=0.5$).\\[2pt]
\textit{Spanish:} imaginary unit = \textit{unidad imaginaria}; complex number = \textit{n\'umero complejo}; complex conjugates = \textit{conjugados complejos}.}}

\wspart{Example 3}{Multiply Complex Numbers}
How can you write each product in the form $a+bi$? \ \textbf{The idea:} distribute as you would with binomials, replace every $i^2$ with $-1$, then write the real part first.

\begin{minipage}[t]{0.48\textwidth}
\textbf{A.} $-2.5i(8-9i)$
\begin{align*}
% TE Example 3A writes this line -2.5i(8)-2.5i(-9i); parentheses added so each product reads as (-2.5i) times a term.
-2.5i(8-9i) &= (-2.5i)(8)+(-2.5i)(-9i)\\
&= -20i+22.5i^2\\
&= -20i+22.5(-1) \quad (i^2=-1)\\
&= -22.5-20i
\end{align*}
The product is $-22.5-20i$ (real part first).

\small\textbf{Watch out:} $22.5i^2$ is $-22.5$, not $22.5$.
\end{minipage}\hfill
\begin{minipage}[t]{0.48\textwidth}
\textbf{B.} $(3-2i)(3+2i)$
\begin{align*}
(3-2i)(3+2i) &= 3(3+2i)-2i(3+2i)\\
&= 9+6i-6i-4i^2\\
&= 9+6i-6i-4(-1)\\
&= 13
\end{align*}
The product is $13$, a real number: $3-2i$ and $3+2i$ are \textbf{complex conjugates} (same real part, opposite imaginary parts), so $6i-6i=0$.
\end{minipage}

\vspace{4pt}
\textbf{Complex conjugates:} $(a+bi)(a-bi)=a^2-abi+abi-b^2i^2=a^2-b^2(-1)=a^2+b^2$, always a real number.

\textbf{Say it:} In Part A, I replaced $i^2$ with \blank[1.5cm] because the \emph{imaginary unit} squared is \blank[1.5cm]. The \emph{real part} of the product is \blank[1.8cm] and the \emph{imaginary part} is \blank[1.8cm]. In Part B, $3-2i$ and $3+2i$ are \emph{complex conjugates}, so the $i$ terms \blank[2.5cm] and the product is \blank[3cm].
\lines[2]

\newpage
\wspart{Try It 3}{Write each product in the form $a+bi$}
\begin{minipage}[t]{0.48\textwidth}
\textbf{a.} $\dfrac{2i}{5}\left(10-\dfrac{5i}{2}\right)$
\lines[6]
\vspace{4pt}
Product: \blank[4cm]
\end{minipage}\hfill
\begin{minipage}[t]{0.48\textwidth}
\textbf{b.} $\left(\dfrac12+2i\right)\left(\dfrac12-2i\right)$
\lines[6]
\vspace{4pt}
Product: \blank[4cm]
\end{minipage}

\vspace{4pt}
In part b, why is there no imaginary part in the answer? Use the words \emph{complex conjugates}.\lines[2]

% Practice (R5): part a is a structural parallel of Savvas quiz item 4 (3i(2-5i)(4+i), a monomial times two
% binomials), which no Try It rehearses; part b is a parallel of quiz item 2 (the bonus: a difference).
\wspart{Practice}{Get ready for quiz questions 4 and 2}
\begin{minipage}[t]{0.48\textwidth}
\textbf{a.} Write the expression $2i(3-4i)(1+i)$ in the form $a+bi$.

\small Multiply the two binomials first (like Part B), then multiply by $2i$ (like Part A).\normalsize
\lines[7]
\vspace{4pt}
Answer: $\blank[1.5cm]+\blank[1.5cm]\,i$
\end{minipage}\hfill
\begin{minipage}[t]{0.48\textwidth}
\textbf{b.} What is the difference of $(4+2i)$ and $(1+7i)$?

\small Distribute the negative to \emph{both} parts first:\\ $(4+2i)+(-1-7i)$.\normalsize
\lines[7]
\vspace{4pt}
Answer: \blank[4cm]
\end{minipage}

\newpage
\wspart{IXL}{Practice for the quiz}
Do the starred skills first. \textbf{SmartScore 70 or above on each starred ($\bigstar$) skill replaces your quiz score.}

\vspace{2pt}
\begin{tabularx}{\textwidth}{l Y Y}
\toprule
\textbf{Code} & \textbf{Skill} & \textbf{Feeds} \\
\midrule
\ixlrow{VZ8 $\bigstar$}{Multiply complex numbers}{Quiz 3, 4: distribute, $i^2=-1$ (Ex.\ 3)}
\ixlrow{7U5 $\bigstar$}{Complex conjugates}{Quiz 3: times its conjugate (Ex.\ 3B)}
\ixlrow{JVF $\bigstar$}{Add and subtract complex numbers}{Quiz 2, the +5 bonus (Practice b)}
\ixlrow{5VV}{Introduction to complex numbers}{Warm-up (optional): $i$ and $a+bi$}
\ixlrow{FG7}{Solve a quadratic equation using square roots}{Quiz 1 (full Savvas quiz only); optional}
\bottomrule
\end{tabularx}

\wspart{Exit question}{Multiply by $i$ (DOK 3)}
The complex number $a+bi$ can be represented on a coordinate plane as the point $(a,b)$. You can use multiplication by $i$ to rotate a point about the origin in the coordinate plane.

\begin{minipage}[t]{0.40\textwidth}
\vspace{0pt}
\begin{tikzpicture}
\begin{axis}[unitgrid, x=0.58cm, y=0.58cm,
  xtick={-5,...,5}, ytick={-5,...,5},
  xticklabels={,$-4$,,$-2$,,,,$2$,,$4$,}, yticklabels={,$-4$,,$-2$,,,,$2$,,$4$,}]
\node[anchor=north east, font=\scriptsize, inner sep=1.5pt] at (axis cs:0,0) {$O$};
\node[closed, red] at (axis cs:3,2) {};
\end{axis}
\end{tikzpicture}
\end{minipage}\hfill
\begin{minipage}[t]{0.57\textwidth}
\vspace{0pt}
\textbf{a.} Write the point $(x,y)$ on the graph as the complex number $x+yi$.\par \blank[4cm]

\vspace{4pt}
\textbf{b.} Multiply the complex number by $i$. Interpret the new value as a new point in the plane.
\lines[2]

\vspace{4pt}
\textbf{c.} Repeat the steps above for two other points. How does multiplication by $i$ rotate a point?
\lines[3]
\end{minipage}

\vspace{6pt}
\textbf{Before the quiz, finish these sentences.}

To multiply complex numbers, I distribute, then replace $i^2$ with \blank[2cm] because $i$ is the \emph{imaginary unit}.

$4+5i$ and $4-5i$ are \emph{complex conjugates}, so their product has no \blank[3.5cm].

In $-22.5-20i$, the \emph{real part} is \blank[2.5cm] and the \emph{imaginary part} is \blank[2.5cm].

\end{document}
